\documentclass[11pt]{article}
\usepackage[margin=1in]{geometry}
\usepackage{amsmath,amssymb,graphicx,booktabs,longtable,array,hyperref,listings,xcolor,fancyhdr}
\hypersetup{colorlinks=true,urlcolor=blue,linkcolor=blue,pdftitle={Regression Diagnostics: Theory, Computation and R and D},pdfauthor={Md. Ikramul Haq Khan}}
\lstset{language=R,basicstyle=\ttfamily\footnotesize,breaklines=true,columns=fullflexible,frame=none,showstringspaces=false}
\setlength{\emergencystretch}{2em}
\setlength{\parindent}{0pt}\setlength{\parskip}{6pt}
\pagestyle{fancy}\fancyhf{}\lhead{Regression diagnostics}\rhead{Md. Ikramul Haq Khan}\cfoot{\thepage}\setlength{\headheight}{14pt}
\title{Model Diagnostics: Industry and Research Practice}
\author{Md. Ikramul Haq Khan}
\date{October 29, 2026}
\begin{document}
\maketitle
\textit{Standalone companion to the complete regression diagnostics reference. Personal study notes prepared with AI assistance; not an official course handout or organizational SOP.}
\section{Start with the scientific question}
Diagnostics assess whether limitations of a fitted model threaten a specific conclusion. Before fitting, record the study design, outcome, observational unit, target population and estimand. Description estimates patterns in observed data; prediction targets future outcomes; causal inference targets effects of interventions and requires identification assumptions beyond residual checks.
For example, predicting industrial stack loss from operating conditions is different from claiming that changing one condition causes an outcome change. A clean residual plot cannot establish absence of confounding.

\section{A practical diagnostic workflow}
\begin{enumerate}
\item Define the target and design. Identify repeated measures, sampling weights, clusters, missingness and temporal order before choosing a model.
\item Check units, coding, impossible values and measurement processes. Preserve legitimate unusual observations.
\item Examine outcome/predictor distributions and scientifically plausible nonlinearities or interactions.
\item Fit a defensible initial model. Inspect residuals versus fitted values, each predictor and relevant time/group variables; inspect a Q--Q plot.
\item Quantify leverage, externally studentized residuals and influence. Investigate flagged cases and compare estimates with and without them as sensitivity analysis, not automatic exclusion.
\item Use targeted tests with stated hypotheses and assumptions. Avoid testing until all p-values exceed 0.05.
\item Compare scientifically justified alternatives. For prediction, use held-out or cross-validated evaluation with splits matching future deployment; keep tuning separate from final evaluation.
\item Report estimates, uncertainty, departures, sensitivity and remaining limitations.
\end{enumerate}
For clustered or survey data, ordinary $H$, Cook's distance and iid reference tests do not directly provide design-valid inference. Use the appropriate weighted/design model and cluster- or PSU-level sensitivity checks. For missing data, complete-case residuals describe the analyzed subset; they do not establish ignorability of missingness. For counts, binary outcomes or infection trajectories, choose appropriate likelihoods and dependence/observation models rather than forcing Gaussian diagnostics onto every outcome.

\section{Industry and major research-group practice: a public-source pathway}
The following translates public resources into training exercises. It is not a claim of access to proprietary pharmaceutical SOPs or internal procedures, nor a regulatory certification. Actual requirements depend on the study, organization and applicable guidance.
\subsection{NIST: measurement, industrial modeling and numerical assurance}
NIST's Statistical Reference Datasets provide certified computational targets. Pair these with the industrial stack-loss lab: separate data quality, numerical accuracy, model adequacy and decision consequences. An R\&D handoff should include the measurement units, calibration/provenance, fitted model, uncertainty, influential operating conditions and limits of extrapolation. Suggested extension (45 minutes): reproduce a linear-regression StRD dataset using QR, verify coefficients and residual SD against the certified values, and document relative/absolute tolerances with explanations.
\subsection{FDA Office of Biostatistics and ICH: clinical-development reasoning}
The FDA's May 2023 covariate-adjustment guidance addresses randomized parallel-group trials and prognostic baseline covariates; its stated scope excludes several other settings, including observational confounding adjustment and longitudinal repeated-measures modeling. Use the document for its actual scope, not as a universal regression manual. ICH E9(R1) supplies an estimand and sensitivity-analysis framework.
Training exercise (60--90 minutes): simulate a two-arm randomized trial with a baseline outcome predictor. Define the population, treatment contrast, outcome, handling of intercurrent events and summary measure before fitting ANCOVA. Prespecify adjustment; compare conventional and justified robust uncertainty, inspect nonlinear baseline relationships and treatment interactions, and assess sensitivity for the same target. Report post-hoc analyses as such. A changed estimand is a different question, not automatically a sensitivity analysis for the original one.
\subsection{Vanderbilt / Frank Harrell: modeling strategy and clinical prediction}
Harrell's publicly available Regression Modeling Strategies material and reproducible-research guidance provide a broader framework: preserve predictor information, consider nonlinear effects, control model complexity, account for missingness and evaluate optimism. A useful diagnostic report connects those choices to performance and uncertainty rather than requiring every plot to look ideal. Exercise (60 minutes): compare prespecified linear and spline forms under resampling; put all adaptive selection inside the resampling loop and evaluate calibration as well as prediction error.
\subsection{WEHI / Bioconductor limma: high-dimensional biomedical R\&D}
The limma documentation, associated with the WEHI bioinformatics resource, illustrates linear modeling for gene-expression experiments and empirical-Bayes variance moderation. The voom method estimates a mean--variance relationship and uses precision weights for RNA-seq analysis. This is a concrete example of adapting the covariance model to measurement technology instead of relying on unweighted Gaussian diagnostics for raw counts.
Exercise (90 minutes or more, optional): use a documented public limma/voom worked example. Check sample metadata and batch structure, inspect sample-level quality and mean--variance trends, verify design rank and estimability, define contrasts and adjust for multiplicity. Do not select genes by unadjusted $p<0.05$ across thousands of tests. Follow the method's own diagnostics; single-outcome OLS residual rules do not transfer unchanged.
\subsection{What a reviewable R\&D diagnostic package contains}
\begin{enumerate}
\item A short scientific question and estimand, population and unit of analysis.
\item A data dictionary, versioned provenance, processing/missingness rules and design information.
\item A prespecified primary model and planned sensitivity checks, with later deviations logged.
\item Stable numerical fitting, rank/conditioning checks and independent checks for critical computations.
\item An evidence table: diagnostic finding, plausible explanation, action, effect on target estimate/uncertainty and remaining uncertainty.
\item Appropriate validation, with participant/site/time splits when those represent intended deployment; no leakage from preprocessing or model selection.
\item Versioned scripts, seeds, package/environment records and an executable report; an independent reviewer can reproduce key outputs.
\end{enumerate}
For an infection-burden thesis, extend this structure to the observation process, reporting delays, temporal dependence and stochastic simulation. A residual distribution check alone cannot validate a transmission mechanism or identify epidemiological parameters.

\section{Additional primary resources for R\&D}
\begin{enumerate}
\item \href{https://www.nist.gov/itl/sed/statistical-reference-datasets}{NIST Statistical Reference Datasets}: numerical benchmarks.
\item \href{https://www.fda.gov/regulatory-information/search-fda-guidance-documents/adjusting-covariates-randomized-clinical-trials-drugs-and-biological-products}{FDA: Adjusting for Covariates in Randomized Clinical Trials (May 2023)}: scope-specific clinical-development guidance.
\item \href{https://www.ema.europa.eu/en/scientific-guidelines/ich-e9-statistical-principles-clinical-trials}{EMA: ICH E9 and E9(R1)}: statistical principles, estimands and sensitivity analysis.
\item \href{https://hbiostat.org/rr/}{Harrell: Reproducible Research}: reproducible statistical collaboration.
\item \href{https://www.bioconductor.org/packages/release/bioc/html/limma}{Bioconductor limma}: current package documentation, user guide and linked primary methods papers. \href{https://bioinf.wehi.edu.au/limma/}{WEHI limma resource}.
\end{enumerate}
\textbf{Maintenance:} Recheck guidance status and installed software documentation before a new project. Retain this dated note as a versioned learning reference, and keep study-specific analysis decisions in the project's own analysis plan.


\end{document}
